\chapter{Teoria Cherna--Weila i dowód wzoru Cherna--Gaussa--Bonneta}
\label{ch:chern-weil}
\index{teoria Cherna--Weila@Teoria Cherna--Weila}

W Sekcji~\ref{sec:gauss-bonnet} całka z Pfaffianu
krzywizny została sformułowana jako wynik
zewnętrzny. Pokażemy, dlaczego krzywizna
reprezentuje klasę Eulera i skąd bierze się
poprawka na brzegu. Ten sam mechanizm określa
klasy Pontriagina potrzebne przy sferach Milnora.

\section{Wielomiany niezmiennicze i transgresja}
\label{sec:chern-weil-26}

Niech $E\to M$ będzie zorientowaną wiązką
rzeczywistą z metryką i koneksją metryczną
$\nabla$. W ramie ortonormalnej macierz
koneksji $\omega$ jest skośnie symetryczna,
a krzywizna ma macierz
$\Omega=\dd\omega+\omega\wedge\omega$.

\begin{definicja}[Wielomian niezmienniczy]
\index{wielomian niezmienniczy@Wielomian niezmienniczy}
\label{def:invariant-polynomial-26}
Symetryczne $r$-liniowe odwzorowanie
$P:\mathfrak{so}(q)^r\to\R$ nazywamy
\emph{niezmienniczym}, gdy
$P(gA_1g^{-1},\ldots,gA_rg^{-1})
=P(A_1,\ldots,A_r)$ dla $g\in\mathrm{SO}(q)$.
Po podstawieniu macierzy form krzywizny
i użyciu iloczynu zewnętrznego otrzymujemy
globalną formę $P(\Omega^r)$ stopnia $2r$.
\end{definicja}

Przykładem jest ślad $P(A,B)=\operatorname{tr}(AB)$.
Dla $q=2m$ polaryzacja Pfaffianu
z~\eqref{eq:gb-pfaff-def} jest
niezmienniczym wielomianem stopnia $m$.
Sprzężenie krzywizny przy zmianie ramki
wyjaśnia, dlaczego te formy sklejają się
z lokalnych wzorów.

\begin{twierdzenie}[Cherna--Weila i transgresja]
\label{thm:chern-weil-26}
Forma $P(\Omega^r)$ jest zamknięta,
a jej klasa de Rhama nie zależy od koneksji.
Dokładniej, dla $\nabla^t=\nabla^0+tA$,
$A=\nabla^1-\nabla^0$, zachodzi
\begin{equation}\label{eq:CW-transgression-26}
 P(\Omega_1^r)-P(\Omega_0^r)
 =\dd\left(r\int_0^1
       P(A,\Omega_t^{r-1})\,\dd t\right).
\end{equation}
Konstrukcja jest naturalna względem cofnięcia wiązki.
\end{twierdzenie}
\begin{proof}
Różnica koneksji $A$ jest globalną jednoformą
endomorfizmową, więc ścieżka $\nabla^t$
ma sens niezależnie od ramki.
Jeśli $D_t$ oznacza różniczkę kowariantną
form endomorfizmowych, to z rozwinięcia
$\Omega_t=\dd\omega_t+\omega_t\wedge\omega_t$
wynika $\frac{\dd}{\dd t}\Omega_t=D_tA$.
Tożsamość Bianchiego daje $D_t\Omega_t=0$.
Niezmienniczość $P$ kasuje wyrazy
komutatorowe z $\omega_t$, więc
\[
 \dd P(\Omega_t^r)=
  rP(D_t\Omega_t,\Omega_t^{r-1})=0,\qquad
 \dd P(A,\Omega_t^{r-1})
  =P(D_tA,\Omega_t^{r-1}).
\]
Symetria $P$ daje jednocześnie
$\frac{\dd}{\dd t}P(\Omega_t^r)
=rP(D_tA,\Omega_t^{r-1})$.
Całkowanie po $t$ kończy dowód.
Równość $F^*\Omega_\nabla=\Omega_{F^*\nabla}$
daje naturalność.
\end{proof}

Forma w nawiasie w~\eqref{eq:CW-transgression-26}
to \emph{forma transgresji} Cherna--Simonsa:
zapisuje zmianę reprezentanta różniczkowego
przy stałej klasie kohomologii.

\begin{przyklad}[Euler i pierwsza klasa Pontriagina]
\label{prz:euler-p1-CW-26}
Dla orientowanej wiązki rzędu $2m$ określamy
$e_\nabla(E)=\operatorname{Pf}(\Omega)/(2\pi)^m$.
Zidentyfikujemy ją z topologiczną klasą
Eulera poniżej. Dla rzeczywistej wiązki
metrycznej normalizacja pierwszej formy
Pontriagina jest
\[
 p_{1,\nabla}(E)=
 -\frac{1}{8\pi^2}\operatorname{tr}
                (\Omega\wedge\Omega).
\]
Jeśli $E$ rozszczepia się na ortogonalną
sumę orientowanych $2$-płaszczyzn,
a koneksja zachowuje ten rozkład
i ma w blokach krzywizny $\omega_i$, blokowy
rachunek śladu daje
$p_{1,\nabla}=\sum_i(\omega_i/2\pi)^2
=\sum_i e_\nabla(E_i)^2$.
Ogólna identyfikacja klasy de Rhama tej
formy z rzeczywistym obrazem całkowitej
klasy $p_1(E)$ korzysta z zasady
rozszczepiania: po cofnięciu na
odpowiednią przestrzeń rozszczepiającą
wiązka dzieli się na takie bloki,
a cofnięcie kohomologii wymiernej
jest injektywne. Topologiczną zasadę
rozszczepiania i wzór $p_1(E_i)=e(E_i)^2$
podaje
\href{https://pi.math.cornell.edu/~hatcher/VBKT/VB.pdf}
{A.\ Hatcher, \emph{Vector Bundles and K-Theory},
§3.1--3.2}; zgodność z formą krzywizny
wynika z powyższego rachunku blokowego
i naturalności, opisanej także przez
\href{https://webusers.imj-prg.fr/~olivier.biquard/m2dg2022.pdf}
{O.\ Biquarda, \emph{Differential and Riemannian Geometry},
§16.c, wzory (III.39) i (III.41)}.
\end{przyklad}

\section{Znormalizowana forma kątowa}
\label{sec:angular-form-26}

Niech teraz $E$ ma rangę $2m$, gdzie $m\geq1$.
Na wiązce sfer jednostkowych
$\pi:S(E)\to M$ istnieje
tautologiczny przekrój $s$ wiązki
$\pi^*E$: w punkcie $(x,u)$
jest nim wektor $u$.
Rozszczepmy
$\pi^*E=\R s\oplus s^\perp$.
Koneksja $\nabla^s$ niech czyni
$s$ równoległym i na $s^\perp$
niech będzie rzutem $\pi^*\nabla$.
Jej forma Eulera jest zerowa,
ponieważ krzywizna ma zerowy
wiersz i kolumnę w kierunku $s$.

\begin{lemat}[Forma kątowa]
\label{lem:angular-form-26}
Niech $\Psi_\nabla$ będzie transgresją
Pfaffianu od $\pi^*\nabla$ do $\nabla^s$:
\[
 \Psi_\nabla=
 \frac{m}{(2\pi)^m}\int_0^1
 P_{\operatorname{Pf}}(A,\Omega_t^{m-1})\,\dd t,
 \qquad A=\nabla^s-\pi^*\nabla.
\]
Wtedy
$\dd\Psi_\nabla=-\pi^*e_\nabla(E)$
oraz $\int_{S(E_x)}\Psi_\nabla=1$
dla każdego włókna zorientowanego
jako brzeg dysku.
\end{lemat}
\begin{proof}
Pierwsza równość jest
\eqref{eq:CW-transgression-26},
bo forma Eulera $\nabla^s$ znika.
Obliczmy normalizację na jednym
włóknie $S^{2m-1}$. W ortogonalnej
trywializacji $\pi^*\nabla$ ogranicza
się do płaskiej koneksji $\dd$.
Niech $u$ będzie wektorem jednostkowym.
Macierz koneksji zachowującej $u$
wynosi
$A=u\,\dd u^{\mathsf T}-\dd u\,u^{\mathsf T}$:
rzeczywiście $Au=-\dd u$.
W punkcie $u=e_1$ mamy
$A_{1a}=\dd u^a$, $A_{a1}=-\dd u^a$,
a $A_{ab}=0$ dla $a,b>1$.
Krzywizna $\dd+tA$ ma wtedy
część styczną
$(\Omega_t)_{ab}=(2t-t^2)
\dd u^a\wedge\dd u^b$:
składnik $\dd A$ daje $2t$,
a $A\wedge A$ daje $-t^2$.
W polaryzacji Pfaffianu jedna
para indeksów zawiera $1$ i
trafia do $A$, zaś pozostałe
parujemy w formach $\Omega_t$.
Po uporządkowaniu iloczynów
wszystkie wyrazy mają ten sam znak,
a ich suma ma współczynnik
$(2m-1)!!$. Stąd
\[
 \Psi_\nabla|_{S^{2m-1}}
 =\frac{(2m-1)!!}{(2\pi)^m}
 \left(\int_0^1(2t-t^2)^{m-1}\dd t\right)
 \operatorname{vol}_{S^{2m-1}}.
\]
Podstawienie $v=1-t$ oraz
całkowanie przez części dają
$\int_0^1(2t-t^2)^{m-1}\dd t
=(2m-2)!!/(2m-1)!!$.
Z rachunku Gaussa
$\operatorname{vol}(S^{2m-1})
=2\pi^m/(m-1)!$.
Ponieważ
$(2m-2)!!=2^{m-1}(m-1)!$,
całka z $\Psi_\nabla$ jest $1$.
\end{proof}

\begin{stwierdzenie}[Identyfikacja klasy Eulera]
\label{prop:euler-identification-26}
Klasa de Rhama $[e_\nabla(E)]$
jest rzeczywistym obrazem
topologicznej klasy Eulera $e(E)$.
\end{stwierdzenie}
\begin{proof}
Na parze $(D(E),S(E))$ forma względna
$(\pi^*e_\nabla(E),-\Psi_\nabla)$
jest zamknięta: w modelu względnego
kompleksu de Rhama warunek brzmi
$\dd\alpha=0$, $i^*\alpha=\dd\beta$,
a tu $\dd(-\Psi_\nabla)=\pi^*e_\nabla(E)$.
Jej całka względna po każdym
$(D^{2m},S^{2m-1})$ wynosi
$0-\int_{S^{2m-1}}(-\Psi_\nabla)=1$.
Jest więc znormalizowaną klasą
Thoma z Twierdzenia~\ref{thm:euler-thom-16}.
Cofnięcie jej przez zerowy przekrój
jest z jednej strony topologiczną
klasą Eulera, a z drugiej klasą
formy $e_\nabla(E)$.
\end{proof}

\section{Dowód i poprawka brzegowa}
\label{sec:CGB-proof-26}

\begin{twierdzenie}[Cherna--Gaussa--Bonneta
z gładkim brzegiem]
\label{thm:CGB-boundary-26}
Niech $(M^{2m},g)$ będzie zwartą,
zorientowaną rozmaitością
riemannowską o gładkim brzegu,
a $\nu$ zewnętrzną jednostkową
normalną. Określmy konkretną
formę brzegową
\begin{equation}\label{eq:CGB-boundary-form-26}
 \mathcal B_g:=\nu^*
       \Psi_{\nabla^{TM}}.
\end{equation}
Wtedy
\begin{equation}\label{eq:CGB-boundary-26}
 \boxed{\displaystyle
 \chi(M)=\frac{1}{(2\pi)^m}
     \int_M\operatorname{Pf}(\Omega^{TM})
       +\int_{\partial M}\mathcal B_g.}
\end{equation}
Jeśli brzeg jest pusty, jest to
Twierdzenie~\ref{thm:chern-gauss-bonnet}.
\end{twierdzenie}
\begin{proof}
Wybierzmy pole wektorowe $V$
równe $\nu$ przy brzegu,
z izolowanymi zerami $p_i$
we wnętrzu. Przedłużamy
$\nu$ przez kołnierz,
a potem perturbujemy pole
we wnętrzu. Usuńmy małe
rozłączne dyski $D_i$
wokół zer i połóżmy
$u=V/|V|$ na pozostałej
części $M_0$. Z Lematu
\ref{lem:angular-form-26}
i twierdzenia Stokesa
otrzymujemy
\[
 \int_{M_0}e_{\nabla^{TM}}
 =-\int_{\partial M_0}u^*\Psi
 =-\int_{\partial M}\nu^*\Psi
   +\sum_i\int_{\partial D_i}u^*\Psi.
\]
Na brzegu dziury znak orientacji
jest przeciwny do orientacji
$\partial D_i$. Gdy dyski
kurczą się do zer, ich
całki zbiegają do stopni
odwzorowań
$u:\partial D_i\to S^{2m-1}$,
czyli indeksów $V$:
część formy na włóknie ma
całkę $1$, a wkład
zmienności lokalnej ramki
zanika wraz z promieniem.
Zatem
\[
 \int_M e_{\nabla^{TM}}
 +\int_{\partial M}\nu^*\Psi
 =\sum_i\operatorname{ind}_{p_i}V.
\]
Suma po prawej wynosi
$\chi(M)$ dla pola
skierowanego na zewnątrz.
Wyprowadźmy używaną
wersję Poincarégo--Hopfa:
wybieramy funkcję Morse'a
stałą na brzegu, rosnącą
w kierunku zewnętrznym.
Jej gradient jest zewnętrzny,
a zero indeksu Morse'a
$\lambda$ ma indeks pola
$(-1)^\lambda$ i wnosi
ten sam składnik do sumy
naprzemiennej uchwytów.
Każde pole zewnętrzne
jest homotopijne na brzegu
przez pola zewnętrzne,
bo ich normalne składowe
są dodatnie; suma indeksów
nie zmienia się przy tej
homotopii. To kończy dowód.
\end{proof}

Dla $m=1$ forma kątowa mierzy
obrót jednostkowego wektora.
Wzdłuż brzegu obrót $\nu$
ma prędkość $k_g$, więc
$\mathcal B_g=k_g\,\dd s/(2\pi)$,
zgodnie z~\eqref{eq:gauss-bonnet-brzeg}.
W ogólnym wymiarze
\eqref{eq:CGB-boundary-form-26}
definiuje formę, która
w~\eqref{eq:chern-gb-brzeg-idea}
była tylko zapowiedziana.
Rozwinięcie transgresji
zawiera drugą formę
fundamentalną brzegu
oraz jego krzywiznę.

\begin{przyklad}[Płaska kula]
\label{prz:CGB-ball-26}
Na $D^{2m}\subset\R^{2m}$
krzywizna znika. Normalna
zewnętrzna $\nu(u)=u$
na $S^{2m-1}$ jest
identycznością włókna,
więc Lemat~\ref{lem:angular-form-26}
daje $\int_{\partial D^{2m}}\mathcal B_g=1
=\chi(D^{2m})$.
Sprawdza to współczynnik
i znak poprawki brzegowej.
\end{przyklad}

\par\smallskip
{\footnotesize\noindent\emph{Dalsza lektura:}
\href{https://webhomes.maths.ed.ac.uk/~v1ranick/papers/chern7.pdf}
{S.-S.\ Chern, \emph{A Simple Intrinsic Proof of the
Gauss--Bonnet Formula for Closed Riemannian Manifolds}
(1944), §1--2},
\href{https://www.imar.ro/~sergium/fisiere/pfa.pdf}
{D.\ Cibotaru i S.\ Moroianu,
\emph{Odd Pfaffian Forms},
część o transgresji i brzegu}.\par}
